\chapter{Walk-forward validation}\label{ch:walkfwd}

\noindent A single holdout split answers ``did it work once?'' Walk-forward validation
answers the harder question: ``does it keep working as the world moves?'' Done naively
it leaks---the label windows at every fold boundary overlap, so the test set remembers
the training set. This chapter is the procedure and its two failure modes.

\section{Anchored vs rolling}\label{sec:wftypes}

Two window conventions. \emph{Anchored}: train on everything up to $t$, test on the next
block, grow the training set each fold. \emph{Rolling}: train on a fixed-width window
that slides with $t$. \texttt{wf\_demo.py} on a three-regime series (trend, quiet,
volatile):
\begin{lstlisting}[numbers=none]
anchored walk-forward:
  train 0:200  test 200:400  IS +2.35  OOS +1.86
  train 0:400  test 400:600  IS +2.11  OOS +0.55
  train 0:600  test 600:800  IS +1.51  OOS -0.52
  train 0:800  test 800:1000  IS +1.16  OOS +1.15
rolling (400-bar) walk-forward:
  train 0:400  test 400:600  IS +2.11  OOS +0.55
  train 200:600  test 600:800  IS +1.01  OOS -0.52
  train 400:800  test 800:1000  IS -0.53  OOS +1.15
  train 600:1000  test 1000:1200  IS -0.88  OOS +0.52
\end{lstlisting}

\begin{example}\label{ex:wf}
\textbf{Reading the fold table.} The same momentum rule: IS decays $2.35\to1.16$ as the
training window dilutes the trending regime; the rolling IS flips sign entirely
($-0.88$) once the trend is out of window.
\begin{solution}
Anchored windows average over history---good when the process is stable, bad when regimes
switch. Rolling windows stay local---good when relationships decay, bad when the window
is too short to estimate anything. The choice is a belief about the data's memory, made
before the data is touched.
\end{solution}
\end{example}

\begin{figure}[htbp]\centering
\begin{tikzpicture}[font=\scriptsize]
  \foreach \y/\a/\b in {3.4/0/200, 2.6/0/400, 1.8/0/600, 1.0/0/800}{
    \fill[accent!30] ({\a/100*1.15},\y) rectangle ({\b/100*1.15},\y+0.55);
    \fill[amber!50] ({\b/100*1.15},\y) rectangle ({\b/100*1.15+2.3},\y+0.55);}
  \node[left] at (-0.1,3.68) {test 200--400};
  \node[left] at (-0.1,2.88) {test 400--600};
  \node[left] at (-0.1,2.08) {test 600--800};
  \node[left] at (-0.1,1.28) {test 800--1000};
  \node[text=accent] at (2.0,4.2) {train (anchored)};
  \node[text=amber!70!black] at (7.0,4.2) {test 200};
  \draw[->,slate] (0,-0.1)--(9.4,-0.1) node[below left=-1pt]{bar index};
\end{tikzpicture}

\caption{Anchored walk-forward: the train window grows; each fold's test block is fresh.}
\label{fig:wf}
\end{figure}

\begin{figure}[htbp]\centering
\begin{tikzpicture}
\begin{axis}[width=0.86\linewidth,height=0.32\linewidth,
    xlabel={bar},ylabel={cum.\ return},
    xmin=0,xmax=1200,ymin=-0.4,ymax=0.7,
    tick style={draw=none},axis line style={draw=slate},font=\scriptsize]
  \addplot[ink,thick] file{figures/data/REG.tex};
  \draw[accent,dashed] (axis cs:400,0.3)--(axis cs:400,0.75);
  \draw[accent,dashed] (axis cs:800,0.3)--(axis cs:800,0.75);
  \node[font=\scriptsize,text=slate] at (axis cs:200,0.68) {trend};
  \node[font=\scriptsize,text=slate] at (axis cs:600,0.68) {quiet};
  \node[font=\scriptsize,text=slate] at (axis cs:1000,0.68) {volatile};
\end{axis}
\end{tikzpicture}

\caption{Why window choice matters: the process itself changes regime (trend, quiet,
volatile). A fixed training span estimates the wrong world in at least one regime.}
\label{fig:regime}
\end{figure}

\section{The boundary leak: purge and embargo}\label{sec:purge}

K-fold CV assumes exchangeability. Financial labels violate it twice: the label is a
\emph{window} (e.g.\ the next-30-bar return), and windows straddling a fold boundary see
both sides. \texttt{purged\_kfold.py} counts the leak honestly:
\begin{lstlisting}[numbers=none]
label window W=30  adjacent-label corr 0.983
boundary 118: 29 train labels reach into test fold
boundary 236: 29 train labels reach into test fold
boundary 354: 29 train labels reach into test fold
boundary 472: 29 train labels reach into test fold
contiguous CV shares 116 overlapping labels across 4 boundaries
purge them + embargo W on each side, or the test is not a test
\end{lstlisting}

\begin{example}\label{ex:purge}
\textbf{What the embargo buys.} A 30-bar label makes adjacent labels $0.98$-correlated;
116 train labels literally extend into the test folds.
\begin{solution}
Purging removes every train row whose label window overlaps the test block; the embargo
adds a margin so features computed over windows do not straddle either. Without both,
``out-of-sample'' scores include a neighbour's answer---the leak is not in the model,
it is in the labelling.
\end{solution}
\end{example}

\begin{figure}[htbp]\centering
\begin{tikzpicture}[font=\scriptsize]
  \draw[->,slate] (0,0)--(9.4,0) node[below left=-1pt]{bar};
  \draw[slate] (0,2.6)--(9.2,2.6) node[right]{price};
  \draw[accent,thick] (1.0,2.9)--(4.0,2.9) node[midway,above]{train label $i$ (30 bars)};
  \draw[amber,thick] (4.3,2.3)--(7.3,2.3) node[midway,below]{test label $j$ (30 bars)};
  \draw[<->,red,thick] (4.0,1.6)--(4.3,1.6) node[midway,below]{overlap};
  \node[align=left,text=slate] at (7.6,3.3) {$\mathrm{corr}(y_i,y_{i+1})=0.983$\\ with a 30-bar label};
\end{tikzpicture}

\caption{A 30-bar forward label: consecutive labels share 29 of 30 windows
($\mathrm{corr}=0.983$).}
\label{fig:overlap}
\end{figure}

\begin{figure}[htbp]\centering
\begin{tikzpicture}[font=\scriptsize]
  \node[align=left] at (0,2.5) {contiguous:};
  \foreach \x/\c in {0/accent!30,2.2/accent!30,4.4/amber!60,6.6/accent!30,8.8/accent!30}{
    \fill[\c] (\x,2.6) rectangle ({\x+2.0},3.1);}
  \draw[<->,red,thick] (4.35,2.4)--(4.45,2.4) node[below,text=red]{leak};
  \node[align=left] at (0,0.9) {purged+embargo:};
  \fill[accent!30] (0,1.0) rectangle (2.0,1.5);
  \fill[accent!30] (2.2,1.0) rectangle (3.9,1.5);
  \fill[white,draw=red,dashed] (3.9,1.0) rectangle (4.7,1.5);
  \fill[amber!60] (4.7,1.0) rectangle (6.7,1.5);
  \fill[white,draw=red,dashed] (6.7,1.0) rectangle (7.5,1.5);
  \fill[accent!30] (7.5,1.0) rectangle (8.9,1.5);
  \fill[accent!30] (9.1,1.0) rectangle (11.0,1.5);
  \node[text=red] at (5.6,0.4) {embargoed margin $W$ on each side};
\end{tikzpicture}

\caption{Contiguous K-fold leaks at the boundary; purged+embargoed folds excise the
overlapping margin.}
\label{fig:purge}
\end{figure}

\section{Per-fold reporting}\label{sec:folds}

The deliverable of walk-forward is not one number---it is the sequence of fold scores and
its dispersion. A strategy that posts $+1.9,+0.6,-0.5,+1.2$ has a different claim than
one posting $+0.9$ four times, even at the same mean.

\begin{figure}[htbp]\centering
\begin{tikzpicture}
\begin{axis}[width=0.86\linewidth,height=0.36\linewidth,
    xlabel={in-sample Sharpe},ylabel={out-of-sample Sharpe},
    xmin=-1,xmax=3,ymin=-1,ymax=2.5,
    tick style={draw=none},axis line style={draw=slate},font=\scriptsize]
  \addplot[only marks,mark=*,mark size=2.5pt,accent] file{figures/data/FOLDS.tex};
  \addplot[slate,dashed] coordinates {(-1,-1)(3,3)};
  \node[font=\scriptsize,text=slate] at (axis cs:2.2,1.4) {OOS = IS line};
\end{axis}
\end{tikzpicture}

\caption{IS vs OOS per fold: points below the diagonal paid the regime tax.}
\label{fig:folds}
\end{figure}

\begin{example}\label{ex:foldread}
\textbf{The fold-reading heuristic.} Fold 3 OOS is $-0.52$ while IS is $+1.51$; fold 4
OOS recovers to $+1.15$.
\begin{solution}
The $-0.52$ fold is the regime break (the quiet stretch)---the strategy still works in
trending and volatile regimes but not in the quiet one. That is information, not failure:
a strategy that is honest about \emph{where} it fails can be gated on regime; a single
pooled number hides exactly the fact you need to size it.
\end{solution}
\end{example}

\section{Putting it together}\label{sec:wfprotocol}

The protocol: pick the window family (anchored if the world is stable, rolling if it
decays); purge label overlaps and embargo $W$ bars on each side of every boundary;
score each fold; report the sequence, the mean, and the worst fold---the worst fold is
the number the risk model should believe.

\begin{figure}[htbp]\centering
\begin{tikzpicture}[font=\scriptsize,every node/.style={draw=slate,fill=white,rounded corners=1pt,align=center}]
  \node (w) {window choice:\\ anchored or rolling};
  \node (p) [right=1.3cm of w] {purge + embargo\\ at every boundary};
  \node (f) [right=1.3cm of p] {score each fold\\ separately};
  \node (d) [right=1.3cm of f] {report the\\ spread};
  \draw[->,ink] (w)--(p); \draw[->,ink] (p)--(f); \draw[->,ink] (f)--(d);
  \node[below=0.5cm of p,draw=none,text=slate] {one number is a fold; the strategy is the sequence};
\end{tikzpicture}

\caption{The walk-forward protocol.}
\label{fig:wfproto}
\end{figure}

\begin{example}\label{ex:nested}
\textbf{Where hyperparameter search fits.} If you tune on the test folds you rebuild
the leak at a higher level.
\begin{solution}
Nest it: inside each walk-forward fold, tune on the train block only (its own inner CV),
freeze, then score once on the test block. The outer fold sequence is then an honest
measure of the \emph{procedure}---signal plus tuner---not of a lucky parameter.
\end{solution}
\end{example}

\begin{checkpoint}[Validation audit]
\begin{itemize}
\item What does the embargo remove that the purge misses? (Feature windows straddling the
boundary, not just label windows.)
\item Which window family suits a decaying relationship? (Rolling.)
\item Why report the worst fold? (It is the number the risk model should believe.)
\end{itemize}
\end{checkpoint}

\section{Recap}\label{sec:wfrecap}

Walk-forward is CV for a world that moves. The boundary leak is mechanical---labels are
windows, windows overlap---and the fix is equally mechanical: purge the overlap, embargo
the margin, and score the sequence. A strategy is the set of its folds, not the average
of them.

\begin{table}[htbp]\centering\small
\caption{Chapter toolbox.}
\begin{tabular}{@{}lll@{}}
\toprule
Concept & Rule & Gotcha\\
\midrule
Anchored & stable worlds & dilutes regime change\\
Rolling & decaying memory & window too short to estimate\\
Purge & remove overlapping labels & labels are windows\\
Embargo & margin $W$ each side & features straddle too\\
Folds & report the sequence & worst fold is the honest number\\
\bottomrule
\end{tabular}
\end{table}

\begin{example}[How many folds is enough?]\label{ex:foldcount}
\textbf{Strategy.} The same walk-forward protocol, sliced into
2, 3, 4 and 6 test folds on identical data. Watch what happens to
the OOS mean and its spread as the count grows.

\textbf{Solution.} \texttt{fold\_count.py} prints:
\begin{verbatim}
2 folds: mean +0.68  sd 0.15     OOS: +0.53 +0.83
3 folds: mean -0.37  sd 0.69     OOS: +0.46 -1.23 -0.33
4 folds: mean +0.43  sd 0.60     OOS: +0.55 -0.52 +1.15 +0.52
6 folds: mean +0.21  sd 0.85
   OOS: +0.88 -0.08 -0.57 +1.65 -0.83 +0.18
more folds -> the estimate fragments; the mean is not the spread
\end{verbatim}
The \emph{same} strategy, the \emph{same} data, and the OOS mean
runs $+0.68 \to -0.37 \to +0.43 \to +0.21$ purely as a function
of how you cut it. Each fold's sd grows with count --- finer
slices give more, noisier, estimates.

\textbf{Significance.} There is no ``correct'' fold count: the
right question is whether the \emph{conclusion} survives across
cuts. Report OOS as mean \emph{and} sd, and re-run at several
counts --- a result that only exists at one slicing is a slicing
artifact.
\end{example}

\section{Common mistakes}
\begin{itemize}
\item \textbf{Shuffled K-fold on time series.} Destroys the arrow of time; the test is
  not a test.
\item \textbf{Purging labels but not features.} A 20-day moving average still crosses
  the boundary.
\item \textbf{One-fold verdicts.} ``It worked on the 2019 split'' is a sample of one.
\item \textbf{Tuning on the outer fold.} Rebuilds the leak at meta level.
\item \textbf{Anchored on a non-stationary world.} The training set remembers a dead regime.
\item \textbf{Averaging away the worst fold.} The dispersion \emph{is} the result.
\item \textbf{Picking the fold count after looking.} OOS
mean runs $+0.68 \to -0.37$ as folds go $2\to3$
(Example~\ref{ex:foldcount}) --- choosing the count that tells
your favourite story is just a subtler leak.
\item \textbf{Averaging without the spread.} Mean OOS Sharpe
$+0.43$ with sd $0.60$ is a coin toss; the sd \emph{is} the
answer, not a footnote.\end{itemize}

\section{Exercises}
\begin{exercise}[easy]
Why does shuffled K-fold fail on time series?
\end{exercise}
\begin{solution}
It leaks the future into training and destroys the causal structure the test is meant
to measure.
\end{solution}

\begin{exercise}[easy]
With a 30-bar label, how many train labels overlap each fold boundary?
\end{exercise}
\begin{solution}
$W-1=29$ per side of the boundary.
\end{solution}

\begin{exercise}[easy]
Anchored vs rolling: which assumes stable relationships?
\end{exercise}
\begin{solution}
Anchored---it weights all history equally; rolling assumes memory decays.
\end{solution}

\begin{exercise}[easy]
What does the embargo remove beyond the purge?
\end{exercise}
\begin{solution}
Train rows whose \emph{feature} windows (not just label windows) cross the boundary.
\end{solution}

\begin{exercise}[easy]
Which fold's score should the risk model trust most, and why?
\end{exercise}
\begin{solution}
The worst fold---it is the realized bad-regime outcome, not an average that hides it.
\end{solution}

\begin{exercise}
A 5-fold walk-forward on 600 rows with $W{=}30$: how many purged rows total?
\end{exercise}
\begin{solution}
$4\times(2\times 29)\approx 232$---the boundaries dominate a small dataset.
\end{solution}

\begin{exercise}
Design nested walk-forward for a two-parameter momentum strategy.
\end{exercise}
\begin{solution}
Inside each outer fold's train block, run an inner walk-forward to pick $(f,s)$; freeze
and score once on the outer test block.
\end{solution}

\begin{exercise}
Why does OOS $=-0.52$ in a quiet regime not falsify the strategy?
\end{exercise}
\begin{solution}
Regime dependence is a property to model (gate or size on regime), not necessarily a
failure of the edge in the regimes where it exists.
\end{solution}

\begin{exercise}
What is the difference between purging and embargoing in one sentence each?
\end{exercise}
\begin{solution}
Purge drops train labels overlapping the test block; embargo drops a safety margin of
$W$ bars on each side so nothing else straddles.
\end{solution}

\begin{exercise}
A colleague reports CV score 0.94 on 30-bar labels. Which audit question first?
\end{exercise}
\begin{solution}
``What was the adjacent-label correlation and were boundary rows purged?''---0.94 on
0.98-correlated labels is partly bookkeeping.
\end{solution}

\begin{exercise}
Why is the rolling window a belief, not a technique?
\end{exercise}
\begin{solution}
Its width encodes an assumed memory timescale of the process---a modelling choice made
before seeing the data.
\end{solution}

\begin{exercise}
Give two reasons the anchored IS decays across folds even for a real edge.
\end{exercise}
\begin{solution}
Regime dilution (the window averages different worlds) and signal crowding/decay over
calendar time.
\end{solution}

\begin{exercise}[hard]
Prove that with a $W$-bar forward label, contiguous K-fold shares at least $W-1$
labels per boundary, and derive the honest minimum embargo.
\end{exercise}
\begin{solution}
Labels $t{-}W{+}1,\ldots,t{-}1$ all contain test bars in their windows: $W{-}1$ rows.
Embargo $\ge W$ clears every feature and label window.
\end{solution}

\begin{exercise}[hard]
Explain why nested walk-forward is the honest way to report a tuner, and what exactly it
measures.
\end{exercise}
\begin{solution}
It treats (signal + selection rule) as the unit under test; the outer sequence measures
how the \emph{procedure} performs OOS, separating the edge from lucky parameters.
\end{solution}

\begin{exercise}[hard]
Propose a walk-forward design for a strategy with monthly labels on a 5-year dataset---
what is the binding constraint?
\end{exercise}
\begin{solution}
Only $\sim$60 labels exist and each is 30-bar-correlated; folds shrink to the point where
per-fold estimates are noise---the binding constraint is effective sample size, not
compute.
\end{solution}

\begin{exercise}[easy]
Why does the OOS mean flip sign between 2 and 3 folds on the same
data?
\end{exercise}
\begin{solution}
Each fold re-weights the regimes: the 3-fold cut puts the quiet
and volatile regimes into separate test windows, so the OOS mean
moves with the composition of the slices --- not with the
strategy.
\end{solution}
\begin{exercise}[medium]
Your walk-forward gives mean OOS $+0.43$, sd $0.60$ at 4 folds.
Compute an approximate 95\% interval for the mean and state what
conclusion it supports.
\end{exercise}
\begin{solution}
$\pm 1.96\cdot0.60/\sqrt4 \approx \pm0.59$: interval
$[-0.16, +1.02]$ includes zero --- the evidence is weak; report
the spread before the point.
\end{solution}
\begin{exercise}[hard]
Argue for or against: ``more folds is always better.'' Use the
bias-variance trade-off on the printed numbers.
\end{exercise}
\begin{solution}
Against. More folds give more estimates (variance of the
\emph{mean} falls like $1/\sqrt{k}$) but each estimate sees a
thinner window and drifts from the regime mix --- at 6 folds the
sd is $0.85$, the largest printed. Fold count trades
\emph{granularity} for \emph{reliability}.
\end{solution}

\section*{Notes}
Validation is where the previous chapter's caution becomes procedure. The leaks here are
not exotic---they are label windows and boundary rows, countable and fixable. What you
cannot fix is a world that changes regime, and that is what the fold sequence is for.

\section*{Further reading}
\begin{itemize}
\item L.~de Prado, \emph{Advances in Financial ML} --- purged CV and embargo (Ch.~7):

\url{https://www.wiley.com/en-us/Advances+in+Financial+Machine+Learning-p-9781119482086}
\item scikit-learn \texttt{TimeSeriesSplit} --- the basic implementation:

\url{https://scikit-learn.org/stable/modules/generated/sklearn.model_selection.TimeSeriesSplit.html}
\item D.~Bailey et al., ``Backtest Overfitting'' (JOIM) --- the context:

\url{https://www.davidhbailey.com/dhbpapers/deflated-sharpe.pdf}
\item Hyndman \& Athanasopoulos, \emph{Forecasting: P\&P} --- time-series CV:

\url{https://otexts.com/fpp3/tscv.html}
\end{itemize}

\begin{rulebox}[Watch instead]
\begin{itemize}
\item ``11. Time Series Analysis II'' --- MIT 18.S096 (83:47). Model evaluation for series;
re-verified 2026-10-11.
\url{https://www.youtube.com/watch?v=_IFUfFuyQlU}
\item ``Python for Data Analysis'' --- PyData/McKinney (10:28). The plumbing; re-verified
2026-10-11.
\url{https://www.youtube.com/watch?v=_T8LGqJtuGc}
\end{itemize}
\end{rulebox}
